% Hausarbeit Jura essay starter — hausarbeit-jura.cls (LPPL-1.3c or later)
% The class is supplied by TeX Live; the legal analysis below is original.
\documentclass[utf8,ngerman,automatter,texgyre,centerpagenumber]{hausarbeit-jura}

\title{Verhältnismäßigkeit und algorithmische Verwaltung}
\subtitle{Eine kurze Untersuchung der Begründungspflicht}
\author[Jordan Blake]{Jordan Blake}
\matrikelnummer{20260418}
\semester[2. Fachsemester]{Wintersemester 2026/27}
\prof{Prof. Dr. Mira Santos}
\date{18. September 2026}

\begin{document}
\maketitle

\chapter*{Einleitung}
Behörden setzen zunehmend regelbasierte Systeme ein, um Anträge zu sortieren,
Risiken zu bewerten oder Prüfungen vorzubereiten. Eine technische
Entscheidungshilfe ersetzt dabei nicht die rechtliche Verantwortung der
Behörde. Diese Hausarbeit untersucht deshalb die These, dass eine
automatisierte Vorprüfung nur dann verhältnismäßig ist, wenn Betroffene die
maßgeblichen Gründe nachvollziehen und wirksam anfechten können.

\chapter{Die Begründung als Verfahrensgarantie}
\section{Transparenz und Nachvollziehbarkeit}
Eine Begründung muss mehr leisten als die Mitteilung, ein System habe einen
bestimmten Wert berechnet. Sie soll den Sachverhalt, die angewandte Regel und
den gedanklichen Weg zur Entscheidung erkennen lassen. Bei einem komplexen
Modell kann dies eine verständliche Beschreibung der verwendeten Merkmale,
der Schwellenwerte und der menschlichen Kontrollschritte erfordern. Die
vollständige Offenlegung des Quellcodes ist dafür nicht zwingend der einzige
oder der beste Weg; entscheidend ist, ob die betroffene Person die
entscheidungserheblichen Gründe prüfen kann.

\section{Geeignetheit und Erforderlichkeit}
Die Vorprüfung ist geeignet, wenn sie einen legitimen Verwaltungszweck
fördert. Sie ist nicht erforderlich, wenn ein weniger eingriffsintensives
Verfahren denselben Zweck mit vergleichbarer Zuverlässigkeit erreicht. Ein
regelmäßiger Stichprobenabgleich mit einer vollständig menschlich geprüften
Kontrollgruppe kann deshalb nicht nur die Qualität messen, sondern auch
zeigen, ob die Automatisierung überhaupt einen zusätzlichen Nutzen bringt.

\chapter{Abwägung und Kontrolle}
Die Abwägung muss die Belastung der betroffenen Person und die Gefahr
systematischer Fehlentscheidungen berücksichtigen. Besonders problematisch
ist ein System, das historische Verwaltungsdaten übernimmt, obwohl diese
Daten frühere Zugangshürden widerspiegeln. Eine hohe Trefferquote kann dann
mit einer ungleichen Fehlerverteilung einhergehen.

Ein verhältnismäßiges Verfahren braucht daher mindestens drei Sicherungen:
\begin{enumerate}
  \item eine dokumentierte Zweck- und Datenbegrenzung,
  \item eine verständliche Mitteilung der entscheidungserheblichen Gründe,
  \item eine erreichbare menschliche Überprüfung mit echter
        Änderungsbefugnis.
\end{enumerate}

\chapter*{Schluss}
Algorithmische Vorprüfung ist nicht schon deshalb unzulässig, weil sie
technisch vermittelt ist. Sie wird aber rechtlich angreifbar, wenn ihre
Begründung nur als undurchsichtiger Score erscheint. Transparenz,
Erforderlichkeit und eine wirksame Kontrolle sind deshalb keine nachträgliche
Dokumentation, sondern Voraussetzungen einer fairen Entscheidung.

\end{document}
